\par\Needspace{6\baselineskip}\medskip
\hypertarget{native-113-article084}{}\pdfbookmark[2]{84.}{article-084}
\begin{center}84.\end{center}
Da zwei nach Betrag und Lage gegebene Radiusvektoren zusammen mit einem Bahnelement die ganze Bahn bestimmen, lässt sich aus diesen Angaben auch die \SourceLeaf{114}\emph{Zeit} bestimmen, in der sich der Himmelskörper von einem Radiusvektor zum anderen bewegt. Dabei muss man die Masse des Körpers entweder vernachlässigen oder wenigstens als bekannt ansehen. Wir bleiben bei der ersten Annahme, auf die sich die zweite leicht zurückführen lässt. Umgekehrt ist damit klar, dass zwei nach Betrag und Lage gegebene Radiusvektoren zusammen mit der Zeit, in der der Himmelskörper das dazwischenliegende Bahnstück durchläuft, die ganze Bahn bestimmen. Diese Aufgabe gehört zu den wichtigsten in der Theorie der Bewegung der Himmelskörper. Sie ist recht schwer zu lösen, da der Ausdruck für die Zeit in Abhängigkeit von den Elementen transzendent und überdies recht kompliziert ist. Umso mehr verdient sie es, mit aller Sorgfalt behandelt zu werden. Die später darzustellende Lösung scheint nichts weiter zu wünschen übrig zu lassen. Wir hoffen daher, dass unsere Leser es begrüßen, wenn wir auch die Lösung vor dem Vergessen bewahren, die wir oft benutzt haben, bevor sich jene spätere Lösung ergab. Bei schwierigeren Aufgaben lohnt es sich immer, mehrere Wege einzuschlagen und eine gute Methode auch dann beizubehalten, wenn man einer besseren den Vorzug gibt. Wir beginnen mit der Darstellung dieser früheren Methode.
